\section{Diffusion Policy、RDT 与世界模型方向}

前面几章主要围绕 Transformer/VLM/VLA，本章看 action 生成路线。Diffusion Policy 把扩散模型用于视觉运动策略学习：不是生成图像，而是生成动作轨迹。它对机器人很自然，因为动作序列本身也可以看成需要逐步去噪和生成的对象。

\begin{figure}[H]
\centering
\includegraphics[width=0.96\textwidth]{ch30-06320.jpg}
\caption{Diffusion Policy：通过动作扩散学习 visuomotor policy。投屏帧，约 01:45:20。}
\end{figure}

\begin{knowledgebox}{读图：扩散模型也可以生成 action}
图中展示了 Diffusion Policy 的训练与推理过程。关键迁移是：扩散模型不只用于图像生成，也可以用于 action trajectory 生成，尤其适合连续动作和多峰轨迹。
\end{knowledgebox}

\subsection{RDT-1B：扩散基础模型用于双手操作}

本节继续追问 action 本身应该如何建模。前面 VLM+Action Head 的路线把动作模块作为高频控制器，但仍需要一个更适合连续轨迹、多峰选择和双手协同的生成机制。RDT-1B 的意义在于把 Diffusion Policy 的直觉放大到基础模型尺度，让“动作生成”不再只是一个小策略网络的附属环节。

RDT-1B 将 Diffusion Policy 扩展到更大的双手操作基础模型。它试图处理 unified action space、不同 action head 和双臂 manipulation。它的意义在于，把小模型式动作扩散推进到更大规模、更接近 foundation model 的方向。

\begin{figure}[H]
\centering
\includegraphics[width=0.96\textwidth]{ch31-06571.jpg}
\caption{RDT-1B：双手操作机器人的扩散基础模型。投屏帧，约 01:49:31。}
\end{figure}

\subsection{Prediction with Action 与 VPP}

Prediction with Action 把视觉预测和动作预测放进联合去噪过程，试图把世界模型和 VLA 连接起来。续作 VPP 强调用 predictive visual representations 形成 generalist robot policy。这里的核心不是只输出 action，而是同时理解未来视觉变化和动作后果。

\begin{figure}[H]
\centering
\includegraphics[width=0.96\textwidth]{ch32-07061.jpg}
\caption{Prediction with Action / VPP：通过联合去噪进行视觉策略学习。投屏帧，约 01:57:41。}
\end{figure}

\begin{knowledgebox}{读图：世界模型进入 VLA}
图中包含视频预测、动作和去噪过程。读这张图时要抓住一点：机器人不只需要知道下一步动作，还需要预测动作会如何改变视觉世界，这就是世界模型方向。
\end{knowledgebox}

\subsection{未来方向：UP-VLA 与在线强化学习}

本节从已有系统转向下一阶段问题。当前 VLA 已经能连接视觉、语言和动作，但还没有稳定解决“理解世界”和“预测动作后果”之间的统一，也没有解决真实环境中持续改进的闭环。因此陈建宇把未来方向压缩为两条线：更统一的模型结构，以及更安全可控的在线学习。

最后，陈建宇列出两个未来方向。UP-VLA 试图统一 understanding 和 prediction，把理解与预测合到 embodied agent 的一个模型中。另一个方向是 online reinforcement learning，用在线反馈继续改进 VLA。难点在于，如果直接对整个 VLM/VLA 做 RL，训练可能不稳定；一种策略是先冻结 VLM，只训练 action head，再逐步放开。

\begin{figure}[H]
\centering
\includegraphics[width=0.96\textwidth]{ch33-07592.jpg}
\caption{iRe-VLA / Online RL：通过在线强化学习改进 Vision-Language-Action Model。投屏帧，约 02:06:32。}
\end{figure}

\begin{importantbox}{未来两条线}
第一条线是统一理解、预测和动作，让模型不仅看懂世界，还能预测动作后果；第二条线是在线强化学习，让机器人在真实或仿真环境反馈中持续改进。
\end{importantbox}

\subsection{本章小结}

Diffusion Policy 和 RDT 把动作生成建模推向更强的连续控制；Prediction with Action/VPP 把世界模型引入 VLA；UP-VLA 和在线 RL 则指向未来的统一模型和持续学习。

